\subsection{All unramified local coordinates and the exact full-to-quotient fibres}

\begin{statement}{E.68}{Proved}{Complete unramified local normal forms, triangular comparison maps and every full lifting fibre}
Retain the actual good-reduction local arithmetic action of E.61,
with $K=\mathbb Q(\zeta_3)$, $v\nmid3$, residue cardinality $q$,
and its specified based Frobenius $F$. Put
\[
 \begin{gathered}
 P=\Pi,\qquad M=P/P'',\qquad W=P/\gamma_3I_Q,\\
 G_r=G/\gamma_{r+1}G,\qquad
 L_r(G)=\gamma_rG/\gamma_{r+1}G\quad(G=P,M,W).
 \end{gathered}
 \tag{E68.1}
\]
These are the actual free, metabelian and single-cusp marked
groups of E.41. Let $F_r(G)$ be the matrix induced by the
actual abelian Frobenius on the integral Hall, Chen or
retained-Hall module $L_r(G)$. Define the finite abelian groups
\[
             \mathcal U_r(G)
               =\operatorname{coker}(F_r(G)-1).
 \tag{E68.2}
\]
After the explicit representative choices described below,
there are homeomorphisms of pointed profinite spaces
\[
 \begin{split}
 \operatorname{Sect}_{\rm unr}(G_r)/G_r
       &\simeq\prod_{j=1}^{r}\mathcal U_j(G),\\
 \operatorname{Sect}_{\rm unr}(G)/G
       &\simeq\prod_{j\ge1}\mathcal U_j(G).
 \end{split}
 \tag{E68.3}
\]
Here an unramified section has tame coordinates $(1,z)$:
its geometric inertia is the identity and its actual
Frobenius image is $zF$. The quotient relation is exactly
$z\mapsto g zF(g)^{-1}$.

For either actual quotient $\pi:P\twoheadrightarrow Q$,
$Q=M$ or $W$, write $R=\ker\pi$ and give $R$ the induced
filtration $R\cap\gamma_rP$. Its finite free graded kernel is
\[
 \begin{split}
 K_r(R)&=(R\cap\gamma_rP)/(R\cap\gamma_{r+1}P)\\
       &=\ker\bigl(L_r(P)\longrightarrow L_r(Q)\bigr),\\
 \mathcal V_r(R)&=\operatorname{coker}(F_r|_{K_r(R)}-1).
 \end{split}
 \tag{E68.4}
\]
The integral maps give short exact sequences of finite groups
\[
 0\longrightarrow\mathcal V_r(R)
      \xrightarrow{\ \iota_r\ }\mathcal U_r(P)
      \xrightarrow{\ \bar\pi_r\ }\mathcal U_r(Q)
      \longrightarrow0 .
 \tag{E68.5}
\]
Every unramified target class has full lifts. For a fixed
target represented by $z_QF$, choose any actual lift
$z_0\in P$ of $z_Q$. Its full geometric-conjugacy fibre is
canonically the following nonabelian cohomology set, and
has the displayed compatible coordinates:
\[
 \begin{split}
 \pi^{-1}([z_QF])
 &\simeq H^1_{\rm cts}\!
       \left(\mathbb Z_3,R,\operatorname{Ad}(z_0)F\right)\\
 &\simeq\prod_{r\ge1}\mathcal V_r(R)
       =\prod_{r\ge4}\mathcal V_r(R).
 \end{split}
 \tag{E68.6}
\]
The second homeomorphism uses Hall lifts and finite Smith
representatives; the first is the actual map $r\mapsto rz_0F$.
The fibre in (E68.6) includes every full local section over
that target: E.61 excludes a ramified full preimage of an
unramified target.

In the product coordinates (E68.3), the actual comparison
is the triangular map
\[
 b_r=\bar\pi_r(a_r)+
                  \varepsilon_r(a_1,\ldots,a_{r-1}).
 \tag{E68.7}
\]
Formula (E68.15) below defines each correction from the
unique lower conjugator, its specified Hall lift and the
genuine nonlinear Frobenius on that lift. Formula (E68.17)
gives every point of every fibre recursively, using exactly
the maps in (E68.5). Thus (E68.6) retains the actual
comparison, including all its correction coordinates.

At $K_v=\mathbb Q_2(\zeta_3)$, the following table evaluates
the first four local factors. An entry in column $e$ is
the multiplicity of $\mathbb Z/3^e$; omitted unit Smith
factors contribute the zero group.
\[
\begin{array}{c|r|rrrrr}
 \text{factor}&\text{lattice rank}&e=1&e=2&e=3&e=4&e=5\\ \hline
 \mathcal U_1(P)=\mathcal U_1(M)=\mathcal U_1(W)
       &10&6&2&0&0&0\\
 \mathcal U_2(P)=\mathcal U_2(M)=\mathcal U_2(W)
       &45&15&3&12&1&0\\
 \mathcal U_3(P)=\mathcal U_3(M)=\mathcal U_3(W)
       &330&24&74&74&24&2\\
 \mathcal U_4(P)&2475&537&435&51&319&3\\
 \mathcal U_4(M)&1485&309&252&15&210&3\\
 \mathcal U_4(W)&2466&532&433&51&319&3\\
 \mathcal V_4(P'')&990&228&183&36&109&0\\
 \mathcal V_4(\gamma_3I_Q)&9&5&2&0&0&0
\end{array}
 \tag{E68.8}
\]
At this place every remaining factor is determined by a
finite integral Hall/Chen calculation modulo
$3^{\,r+b_r+1}$, where
$b_r=\max_{r\le j\le2r}(1+v_3(j))$.
This is a proved precision bound for the kernel factors
as well as the three quotient factors.

\emph{Manuscript: \S51.16, folios733v--736v, scans172--178,
formulas(371)--(379), Part E(c),(d). Together with E.66,
this describes every fibre of the local full-to-metabelian
map. For the marked map it describes all unramified fibres;
E.67 retains its separate exact ramified marked-cusp locus.}
\end{statement}

\begin{proof}
\emph{Integral strictness and the actual kernel filtration.}
The free pro-three lower-central graded algebra is the
free integral Lie algebra on $V=P^{\rm ab}$, with its
Hall basis. E.41 proves, through the completed Fox
resolution and the actual commutator map, that
$\operatorname{gr}(M)$ is the integral free metabelian
Lie algebra. In particular its graded pieces are free
and the map from the full Lie algebra is the natural
commutator-to-syzygy map, with kernel $\operatorname{Lie}(V)''$.
For the marked quotient, E.41 proves that the initial
ideal of the actual closed subgroup $\gamma_3I_Q$ is
$\gamma_3\mathfrak i_Q$, where $\mathfrak i_Q$ is the
Lie ideal generated by the actual cusp vector. Its
multihomogeneous Hall basis consists of the words with
at least three occurrences of that vector. The Magnus
inclusion and the lifting of those bracket generators
prove both directions of this equality over
$\mathbb Z_3$, including passage to the closed groups.

Consequently the exact degree-$r$ sequence is
\[
 0\longrightarrow K_r(R)\longrightarrow L_r(P)
       \xrightarrow{\pi_r}L_r(Q)\longrightarrow0,
 \qquad L_r(Q)\text{ finite free over }\mathbb Z_3 .
 \tag{E68.9}
\]
For clarity, the group calculation behind its first term is
\[
 \frac{(R\gamma_{r+1}P)\cap\gamma_rP}{\gamma_{r+1}P}
 =\frac{(R\cap\gamma_rP)\gamma_{r+1}P}{\gamma_{r+1}P}
 \simeq\frac{R\cap\gamma_rP}{R\cap\gamma_{r+1}P}.
\]
The products are closed products of compact subgroups.
Thus the kernel is saturated and is exactly the induced
filtration on $R$, not its intrinsic lower-central
filtration. Every quotient in (E68.9) is torsion-free.
At degrees below four the kernels vanish. At degree four,
$K_4(P'')$ has the Hall basis of brackets of two
degree-two Hall words and rank $\binom{45}{2}=990$;
$K_4(\gamma_3I_Q)$ has basis
$[d,[d,[d,b_j]]]$, $1\le j\le9$, and rank nine.

All three groups equal the inverse limits of the indicated
lower-central quotients, by the Magnus and completed
Fox descriptions in E.41. Since $R$ is a closed subgroup
of $P$, it likewise equals
$\varprojlim R/(R\cap\gamma_{r+1}P)$.
These statements supply the integral groups and
topologies to which the following central induction applies.

\emph{Weights, unique conjugators and unramified sections.}
The local action is the genuine good-reduction action of
E.61 and factors through its pro-three arithmetic quotient.
An arbitrary $z\in G$ defines an unramified section:
the closure of $zF$ in $G\rtimes\mathbb Z_3$ is
procyclic, and projection sends its specified generator
to the specified generator of $\mathbb Z_3$.
The corresponding continuous homomorphism from
$\mathbb Z_3$ is therefore a section. This proves the
actual parametrization by $z$, before conjugacy.

On $V$ the finite-field weights are the rank-two compact
elliptic weights and the cusp Tate weights. On every
degree-$r$ Lie piece their products have complex absolute
values $q^{j/2}$ with $r\le j\le2r$. They survive as
subquotients on all the lattices in (E68.9), so none
has eigenvalue one. In particular every $F_r-1$ in
(E68.2) and (E68.4) is rationally invertible. Its integral
cokernel is finite.

Conjugating the Frobenius action by an arbitrary element
of $P$ changes no associated-graded matrix. A nonidentity
fixed element in any of the finite nilpotent quotients,
or in their separated inverse limits, would have a first
nonzero graded coordinate fixed by $F_r$. Hence
\[
 H^0\bigl(\mathbb Z_3,G,\operatorname{Ad}(z)F\bigr)=1,
 \quad
 H^0\bigl(\mathbb Z_3,R,\operatorname{Ad}(z_0)F\bigr)=1,
 \tag{E68.10}
\]
with the same conclusion at every filtered finite stage.
For $R$ use the actual induced graded lattice $K_r(R)$.
Thus any conjugator between two equivalent Frobenius
representatives is unique.

\emph{Actual recursive representatives.}
Choose ordered integral Hall lifts, or the Chen and
retained-Hall lifts of E.41, giving continuous set sections
$s_r^G:G_{r-1}\to G_r$ with zero last coordinates.
Central induction on the free integral graded bases
proves existence and uniqueness of these collected
coordinates; it uses group commutators and three-adic
powers, without a rational exponential.
Choose a Smith representative map
$\sigma_r^G:\mathcal U_r(G)\to L_r(G)$, with
$\sigma_r^G(0)=0$.
The representative associated to a finite coordinate
sequence is defined literally by
\[
 z_0=1,\qquad
 z_r=s_r^G(z_{r-1})\,\sigma_r^G(a_r).
 \tag{E68.11}
\]
The last factor is viewed in the central kernel of
$G_r\to G_{r-1}$.

To normalize an arbitrary $z\in G_r$, first use the
unique lower conjugator and lift it by $s_r^G$.
The resulting element differs from the chosen lift
of the lower representative by a unique central
coordinate $e\in L_r(G)$.
Its class is $a_r=[e]\in\mathcal U_r(G)$.
There is a unique integral solution
\[
       v=(F_r(G)-1)^{-1}
                      (e-\sigma_r^G(a_r))\in L_r(G).
 \tag{E68.12}
\]
The displayed inverse is only applied to an element
of its integral image, as certified by the Smith
coordinates. Conjugation by the central element $v$
subtracts $(F_r-1)v$, so supplies exactly (E68.11).
This constructs the unique full conjugator and proves
both surjectivity and injectivity of the finite-stage
normal form.

All normalizations are continuous: Smith reduction
has finite residue targets, and the inverse of an
injective lattice map is continuous on its compact
image. The representatives (E68.11) are compatible.
If two inverse-limit sections are equivalent at each
finite stage, the unique conjugators are compatible
and give one element of $G$. Conversely a compatible
sequence of finite normal forms has the actual
representative $\varprojlim z_r$. This proves (E68.3)
as a homeomorphism. The same argument, using the
induced filtration on $R$, gives the second map in
(E68.6).

\emph{The exact fibre as a twisted kernel section set.}
Fix $z_Q$ and $z_0$ as in (E68.6).
For a source class over $[z_QF]$, its quotient
representative is carried to $z_Q$ by a unique
element of $Q$, by (E68.10). Lift that element to
$P$ and apply the actual twisted conjugation.
The resulting representative has the form $rz_0$
with $r\in R$. Any two choices of the lifted
conjugator differ by an element of $R$. Conversely
a conjugacy between two such representatives
projects to a stabilizer of $z_QF$, hence projects
to the identity. Their exact equivalence relation is
\[
 r\longmapsto u r\,(\operatorname{Ad}(z_0)F)(u)^{-1}
                   \qquad(u\in R).
 \tag{E68.13}
\]
This proves the first homeomorphism in (E68.6):
it is a continuous bijection from a compact orbit
space to the Hausdorff fibre.
Every target has a lift because $P\to Q$ is onto.
Changing the chosen lift to $z_0'=r_0z_0$ changes
the representative by the explicit map
$r\mapsto rr_0^{-1}$ and changes the action to
$\operatorname{Ad}(r_0)\operatorname{Ad}(z_0)F$.
Thus all choices are related by the actual inner
cocycle transport.

Both kernels $R$ lie in $P'$, and both quotient maps
induce the identity on $V$. A full ramified section
has nonzero cusp inertia in $V$ by E.61. It therefore
cannot map to an unramified target. This proves that
(E68.6) is the whole full section fibre, not merely
its intersection with a chosen subclass.

\emph{The finite comparison sequence and its nonlinear correction.}
Apply $F_r-1$ to (E68.9). Its three kernels are
zero, so the snake lemma gives exactly (E68.5).
There is also a direct integral inverse description.
Choose any lattice section $s:L_r(Q)\to L_r(P)$.
For $x\in L_r(P)$ whose quotient lies in
$(F_r(Q)-1)L_r(Q)$, let
$y=(F_r(Q)-1)^{-1}\pi_r(x)$.
Then
\[
       x-(F_r(P)-1)s(y)\in K_r(R)
 \tag{E68.14}
\]
represents its unique class in $\mathcal V_r(R)$.
Changing $x$ by a full coboundary changes this
element by a kernel coboundary. This specifies
every map and the inverse on the kernel in (E68.5).

To compute the actual section map, take a source
coordinate prefix and its representative
$z_{r-1}^P$. Let $z_{r-1}^Q$ be the already computed
target normal form. Let $h_{r-1}\in Q_{r-1}$ be
the unique element aligning their quotient
representatives. Put $\widehat h=s_r^Q(h_{r-1})$
and define $e_r\in L_r(Q)$ by the exact group
identity
\[
 \widehat h\,\pi\bigl(s_r^P(z_{r-1}^P)\bigr)
       F(\widehat h)^{-1}
          =s_r^Q(z_{r-1}^Q)\,e_r,
 \qquad \varepsilon_r=[e_r].
 \tag{E68.15}
\]
Both sides have the same lower image, so this
central element exists uniquely. In particular
$F(\widehat h)$ is the genuine Frobenius value,
including its higher coordinates.
Centrality then gives (E68.7). To continue the
comparison, set
\[
 \begin{split}
 v_r&=(F_r(Q)-1)^{-1}
    \bigl(e_r+\pi_r\sigma_r^P(a_r)
                    -\sigma_r^Q(b_r)\bigr),\\
 h_r&=v_r\,\widehat h .
 \end{split}
 \tag{E68.16}
\]
The cokernel equality makes $v_r$ integral.
This is the unique actual aligning conjugator at
the new stage. Equations (E68.15)--(E68.16)
retain all nonlinear corrections and their signs.

Choose for each surjection in (E68.5) a section
of finite sets $\tau_r$ on the specified Smith
residues. For a target coordinate sequence
$(b_r)$, every source in its fibre is obtained
uniquely by the recursion
\[
       a_r=\tau_r\bigl(b_r-\varepsilon_r(a_{<r})\bigr)
                          +\iota_r(k_r),
             \qquad k_r\in\mathcal V_r(R).
 \tag{E68.17}
\]
All maps at any finite stage are maps of finite
sets, so this recursion and its inverse are
continuous on the inverse products.
It is the comparison-coordinate version of the
twisted-kernel homeomorphism (E68.6), with every
chosen representative and correction specified.
No homomorphic splitting of (E68.5) is required.

\emph{Finite exact calculations at the place above two.}
The actual rank-ten matrix in E.37 is
\[
 A_F=\begin{pmatrix}-2I_2&0\\-2U&4I_8\end{pmatrix},
 \qquad U_{d,\bullet}=(v_d,-u_d).
 \tag{E68.18}
\]
The two cusp rows $A_2,A_0$ give a unit minor of
$U$. Integral row operations make them $I_2$
and clear the remaining rows. After permuting
coordinates, (E68.18) is integrally conjugate to
two blocks
$\left(\begin{smallmatrix}-2&0\\-2&4\end{smallmatrix}\right)$
and six scalar blocks $4$. This change fixes the
actual marked vector $e_{B_0}$ and has an integral
inverse, both retained in the certificate.
An actual compatible free pro-three basis is obtained
by taking the ordered products of the original
generators with exponent columns of that inverse
matrix. Their images modulo the Frattini subgroup
are a basis, so they give a free basis; the marked
generator itself is unchanged. All nonlinear
Frobenius values in these new words are retained
in (E68.15)--(E68.16).
Hall multidegrees in these eight invariant
blocks are preserved by Frobenius. Through
degree four their largest block has size24.
The program computes every full, quotient and
kernel matrix in these blocks and saves both
Smith transformations and their inverses.

At degree four the metabelian kernel consists
of the Hall brackets of pairs of degree-two
words, so is exactly
$\bigwedge^2(\bigwedge^2V)$ with its integral
basis. The marked kernel is the nine-word
basis in the proof of (E68.9); its action is
$4^3F_B$. These supply independent checks
on the kernel maps. The determinant valuations
of the full and the two quotient degree-four
operators are $2851,1713,2842$; the kernel
valuations are $1138,9$. Their differences
agree with every exact sequence in (E68.5).
The actual factors are those in (E68.8).

For completeness, the precision bound follows
on all the lattices, not only the finite
computed degrees. The split sublattice
$T\oplus jT_3E$ contains $3V$.
Its degree-$r$ functorial maps $B_r,C_r$
satisfy $B_rC_r=3^rI$, as in E.41.
They preserve the metabelian kernel and the
marked ideal, so restrict to the kernels in
(E68.9) as well as descend to the quotients.
In the split lattice every multihomogeneous
piece has Frobenius scalar $(-2)^j$,
$r\le j\le2r$. The valuations
$v_3((-2)^j-1)=1+v_3(j)$ imply
\[
 3^{\,r+b_r}(F_r-1)^{-1}L_r\subset L_r
 \quad\text{and}\quad
 3^{\,r+b_r}(F_r|_{K_r}-1)^{-1}K_r\subset K_r .
 \tag{E68.19}
\]
Thus all finite factors are killed by that
power, and the stated extra digit suffices
to determine their Smith invariants.

Finally the program checks the nonlinear
normalization separately in the actual
rank-two arithmetic quotient of E.54, where
the Hall coordinates are $(a,b,n,r,s)$ and
\[
 F(a,b,n,r,s)=
  (4a,4b,4b+16n,40n+64r,
       24n+64s-8b(b-1)-15b/2).
 \tag{E68.20}
\]
Its unique normal representatives have
$a,b,n\in\{0,1,2\}$ and $r,s\in\{0,\ldots,8\}$.
Five hundred exact tests recover the complete
conjugator, including nonintegral rational
coordinates with three-adic unit denominators.
For the separate actual class-two projection test,
deliberately choose the alternative nonpointed
target degree-one Hall lift $x^{a+3}y^{b+3}$.
The main normal forms (E68.3) retain their pointed
Smith representatives. In this illustrative test
the identity source has target coordinates
$(0,0,1)$.
The resulting correction in (E68.15) is
$\varepsilon_2(a,b)=1+a-b\pmod3$.
All2187 source representatives give exactly
27 target classes with81 representatives
in each fibre. This verifies the actual
nonlinear correction alongside the952
exact graded sequence certificates; it
does not substitute this smaller quotient
for the rank-ten matrices of (E68.8).

All completed-group and section maps in
the theorem are therefore supplied by the
actual quotient, the unique conjugators,
and the displayed integral coordinate
procedures.
\end{proof}
